% ---------------------------------------------------------------
% preamble/macros.tex - math shortcuts, highlights and notes (every layout)
% ---------------------------------------------------------------

% ===== Number sets =====
\newcommand{\N}{\mathbb{N}}
\newcommand{\Z}{\mathbb{Z}}
\newcommand{\Q}{\mathbb{Q}}
\newcommand{\R}{\mathbb{R}}
\newcommand{\C}{\mathbb{C}}
\newcommand{\F}{\mathbb{F}}
\newcommand{\K}{\mathbb{K}}

% ===== Calligraphic letters: \cA ... \cZ =====
\makeatletter
\@for\@tmpl:=A,B,C,D,E,F,G,H,I,J,K,L,M,N,O,P,Q,R,S,T,U,V,W,X,Y,Z\do{%
  \expandafter\xdef\csname c\@tmpl\endcsname{\noexpand\mathcal{\@tmpl}}}
\makeatother

% ===== Delimiters that size themselves =====
% \abs{x}   fixed size      \abs*{x}   auto \left \right     \abs[\big]{x}   manual
\DeclarePairedDelimiter{\abs}{\lvert}{\rvert}
\DeclarePairedDelimiter{\norm}{\lVert}{\rVert}
\DeclarePairedDelimiter{\floor}{\lfloor}{\rfloor}
\DeclarePairedDelimiter{\ceil}{\lceil}{\rceil}
\DeclarePairedDelimiter{\paren}{(}{)}
\DeclarePairedDelimiter{\brac}{[}{]}
\DeclarePairedDelimiter{\set}{\{}{\}}
\DeclarePairedDelimiterX{\inner}[2]{\langle}{\rangle}{#1,\,#2}
% \setof{x \in \R}{x > 0}  ->  { x in R | x > 0 }
\DeclarePairedDelimiterX{\setof}[2]{\{}{\}}{#1 \;\delimsize\vert\; #2}

% ===== Operators =====
\DeclareMathOperator{\sgn}{sgn}
\DeclareMathOperator{\tr}{tr}
\DeclareMathOperator{\rank}{rank}
\DeclareMathOperator{\diag}{diag}
\DeclareMathOperator{\spn}{span}
\DeclareMathOperator{\img}{im}
\DeclareMathOperator{\Hom}{Hom}
\DeclareMathOperator{\End}{End}
\DeclareMathOperator{\Aut}{Aut}
\DeclareMathOperator{\id}{id}
\DeclareMathOperator{\dom}{dom}
\DeclareMathOperator{\supp}{supp}
\DeclareMathOperator{\lcm}{lcm}
\DeclareMathOperator*{\argmax}{arg\,max}
\DeclareMathOperator*{\argmin}{arg\,min}

% ===== Probability and statistics =====
\DeclareMathOperator{\E}{\mathbb{E}}
\DeclareMathOperator{\Prob}{\mathbb{P}}
\DeclareMathOperator{\Var}{Var}
\DeclareMathOperator{\Cov}{Cov}
\DeclareMathOperator{\Corr}{Corr}
\newcommand{\Normal}{\mathcal{N}}
\newcommand{\iid}{\overset{\text{iid}}{\sim}}

% ===== Calculus =====
% \dd x              upright d with correct spacing: \int f(x) \dd x
% \dv{f}{x}          df/dx         \dv[2]{f}{x}    d^2f/dx^2
% \pdv{f}{x}         partial       \pdv[2]{f}{x}
% \eval{F(x)}{a}{b}  F(x) |_a^b
\newcommand{\dd}{\mathop{}\!\mathrm{d}}
\newcommand{\dv}[3][]{\frac{\dd^{#1} #2}{\dd #3^{#1}}}
\newcommand{\pdv}[3][]{\frac{\partial^{#1} #2}{\partial #3^{#1}}}
\newcommand{\eval}[3]{\left. #1 \right\rvert_{#2}^{#3}}
\newcommand{\limn}{\lim_{n \to \infty}}
\newcommand{\limx}[1]{\lim_{x \to #1}}
\newcommand{\sumn}[1][1]{\sum_{n=#1}^{\infty}}

% ===== Linear algebra =====
% \vb{v} bold vector   \mat{1 & 2 \\ 3 & 4}   \dmat{...} determinant
% \cvec{1 \\ 2 \\ 3} column vector            \vect{a}{b} 2D column vector
\newcommand{\vb}[1]{\symbf{#1}}
\newcommand{\mat}[1]{\begin{bmatrix} #1 \end{bmatrix}}
\newcommand{\dmat}[1]{\begin{vmatrix} #1 \end{vmatrix}}
\newcommand{\cvec}[1]{\begin{bmatrix} #1 \end{bmatrix}}
\newcommand{\vect}[2]{\begin{bmatrix} #1 \\ #2 \end{bmatrix}}
\newcommand{\T}{^{\mathsf{T}}}
\newcommand{\inv}{^{-1}}

% ===== Logic, definitions, misc =====
\newcommand{\defeq}{\coloneqq}          % x := y
\newcommand{\eqdef}{\eqqcolon}          % y =: x
\newcommand{\st}{\txtst}                % "such that" / "โดยที่"
\newcommand{\eps}{\varepsilon}
\newcommand{\vphi}{\varphi}
\newcommand{\bigO}{\mathcal{O}}
\newcommand{\degree}{^{\circ}}
\newcommand{\divides}{\mid}
\newcommand{\ndivides}{\nmid}
\newcommand{\ol}[1]{\overline{#1}}
\newcommand{\wh}[1]{\widehat{#1}}
\newcommand{\wt}[1]{\widetilde{#1}}

% ===== Research notes: visible only in  make draft =====
% \todo{check the constant}   \idea{try induction}   \changed{new text}
\newif\ifdraftmode
\draftmodefalse
\newcommand{\todo}[1]{\ifdraftmode{\color{red!70!black}\textbf{[TODO:} #1\textbf{]}}\fi}
\newcommand{\idea}[1]{\ifdraftmode{\color{blue!70!black}\textbf{[IDEA:} #1\textbf{]}}\fi}
\newcommand{\changed}[1]{\ifdraftmode{\color{green!45!black}#1}\else#1\fi}

% ===== Highlights for summaries =====
% \hl{short phrase}        yellow background, one line only (no wrap)
% \key{important words}    bold red text, wraps normally
% \begin{important} ... \end{important}   whole sentence/paragraph, yellow box
% soul's \hl is not used: it breaks Thai text under XeLaTeX.
\colorlet{hlcolor}{yellow!45}
\colorlet{keycolor}{red!65!black}
\newcommand{\hl}[1]{{\setlength{\fboxsep}{1pt}\colorbox{hlcolor}{\strut #1}}}
\newcommand{\key}[1]{{\bfseries\color{keycolor}#1}}
\newtcolorbox{important}{
  breakable, enhanced, frame hidden, sharp corners,
  colback = hlcolor,
  borderline west = {2pt}{0pt}{orange!80!black},
  left=6pt, right=6pt, top=3pt, bottom=3pt,
  before skip=0.6em, after skip=0.6em,
}

% ===== Formula / key-point box =====
% \begin{formulabox}{Title} ... \end{formulabox}
\newtcolorbox{formulabox}[1]{
  breakable, enhanced,
  colback = blue!3,
  colframe= blue!40!black,
  boxrule = 0.6pt,
  title   = \bfseries #1,
  fonttitle = \normalsize,
}

% ===== Eye-catching note =====
% \note{Do not forget to check the assumptions}
% Orange box with a warning sign, titled "Note" / "หมายเหตุ"
\newcommand{\warnsign}{%
  \tikz[baseline=-0.6ex]{%
    \fill[orange!90!black, rounded corners=0.6pt]
      (0,0.75em) -- (0.48em,-0.1em) -- (-0.48em,-0.1em) -- cycle;
    \node[white, font=\bfseries\footnotesize, inner sep=0pt] at (0,0.24em) {!};}}
\newcommand{\note}[1]{%
  \begin{tcolorbox}[
    breakable, enhanced, frame hidden, sharp corners,
    colback = orange!8,
    borderline west = {2.5pt}{0pt}{orange!90!black},
    left=6pt, right=6pt, top=4pt, bottom=4pt,
    before skip=0.6em, after skip=0.6em,
  ]
    {\large\warnsign}\ {\bfseries\color{orange!60!black}\txtnote}\quad #1
  \end{tcolorbox}}

% ===== Page x of y =====
% \pageofpages  ->  "หน้า 2/5"
\newcommand{\pageofpages}{\txtpage~\thepage/\pageref*{LastPage}}

% ===== Thai inside \texttt and code =====
% The mono font has no Thai glyphs, so each run of Thai characters in the
% tt family switches to \thaifallbackfont and back. Other text is untouched.
\makeatletter
\newXeTeXintercharclass\thai@class
\@tempcnta="0E01
\loop\ifnum\@tempcnta<"0E5C
  \XeTeXcharclass\@tempcnta=\thai@class
  \advance\@tempcnta by 1
\repeat
\newif\ifthai@inmono
\newcommand{\thai@enter}{%
  \edef\thai@cur{\f@family}\edef\thai@tt{\ttdefault}%
  \ifx\thai@cur\thai@tt\thai@inmonotrue\thaifallbackfont\fi}
\newcommand{\thai@leave}{\ifthai@inmono\thai@inmonofalse\ttfamily\fi}
\XeTeXinterchartoks 0 \thai@class = {\thai@enter}
\XeTeXinterchartoks 4095 \thai@class = {\thai@enter}
\XeTeXinterchartoks \thai@class 0 = {\thai@leave}
\XeTeXinterchartoks \thai@class 4095 = {\thai@leave}
\XeTeXinterchartokenstate=1
\makeatother
